\Chapter{55}

\Verse{1}
{
    \zA{且说荣府中刚将年事忙过，凤姐儿因年内年外操劳太过，一时不及检点，便小 月了，不能理事，天天两三个大夫用药。 凤姐儿自恃强壮，虽不出门，然筹画计算，
        想起什么事来，就叫平儿去回王夫人，任人谏劝，他只不听。 王夫人便觉失了膀臂， 一人能有多少精神？ 凡有了大事就自己主张，将家中琐碎之事，一应都暂令李纨协
        理。 李纨本是个尚德不尚才的，未免逞纵了下人，王夫人便命探春合同李纨裁处， 只说过了一月，凤姐将养好了，仍交给他。 谁知凤姐？}
}
{
    \eA{We will now resume our narration with the Jung Mansion. Soon after the
        bustle of the new year festivities, lady Feng who, with the most arduous
        duties she had had to fulfil both before and after the new year, had found
        little time to take proper care of herself, got a miscarriage and could not
        attend to the management of domestic affairs. Day after day two and three
        doctors came and prescribed for her.}
}
{
}

\Verse{2}
{
    \zA{赋气血不足，兼年幼不知保 养，平生争强斗智，心力更亏，故虽系小月，竟着实亏虚下来。 一月之后，又添了 下红之症。
        他虽不肯说出来，众人看他面目黄瘦，便知失于调养。 王夫人只令他好 生服药调养，不令他操心。 他自己也怕成了大症，遗笑于人，便想偷空调养，恨不
        得一时复旧如常。 谁知服药调养，直到三月间，才渐渐的起复过来，下红也渐渐止 了。 此是后话。}
}
{
    \eA{But lady Feng had ever accustomed herself to be hardy, so although unable to
        go out of doors, she nevertheless devised the ways and means for everything,
        and made the various arrangements she deemed necessary, and whatever concern
        suggested itself to her mind, she entrusted to P'ing Erh to lay before
        Madame Wang. But however much people advised her to be careful, she would
        not lend an ear to them.}
}
{
}

\Verse{3}
{
    \zA{如今且说目今王夫人见他如此，探春和李纨暂难谢事，园中人多，又恐失于照 管，特请了宝钗来，托他各处小心。 因嘱咐他：“老婆子们不中用，得空儿吃酒斗
        牌，白日里睡觉，夜里斗牌，我都知道的。 凤丫头在外头，他们还有个怕惧，如今 他们又该取便了。 好孩子，你还是个妥当人，你兄弟妹妹们又小，我又没工夫，你
        替我辛苦两天照应照应。 凡有想不到的事你来告诉我，别等老太太问出来我没话 回。 那些人不好你只管说，他们不听你来回我。 别弄出大事来才好。”}
}
{
    \eA{Madame Wang felt as if she had been deprived of her right arm. And as she
        alone had not sufficient energy to see to everything, she bestowed her own
        attention upon such important affairs, as turned up, and entrusted, for the
        time being, all miscellaneous domestic matters to the co-operation of Li
        Wan.}
}
{
}

\Verse{4}
{
    \zA{宝钗听说， 只得答应了。 时届季春，黛玉又犯了咳嗽； 湘云又因时气所感，也病卧在蘅芜院，一天医药 不断。
        探春和李纨相住间壁，二人近日同事，不比往年，往来回话人等亦甚不便， 故二人议定，每日早晨，皆到园门口南边的三间小花厅上去会齐办事，吃过早饭，
        于午错方回。 这三间厅原系预备省亲之时众执事太监起坐之处，故省亲以后也用不 着了，每日只有婆子们上夜。 如今天已和暖，不用十分修理，只不过略略的陈设些，
        便可他二人起坐。}
}
{
    \eA{Li Wan had at all times held virtue at a high price, and set but little
        value on talents of any kind, so that she, as a matter of course, displayed
        leniency to those who were placed under her. Madame Wang accordingly bade
        T'an Ch'un combine with Li Wan in the management of the household. "In a
        month," she argued, "lady Feng will be getting all right again, and then you
        can once more hand over charge to her."}
}
{
}

\Verse{5}
{
    \zA{这厅上也有一处匾，题着“辅仁谕德”四字，家下俗语皆只叫“议 事厅儿”。 如今他二人每日卯正至此，午正方散，凡一应执事的媳妇等来往回话的， 络绎不绝。
        众人先听见李纨独办，各各心中暗喜，因为李纨素日是个厚道多恩无罚的人， 自然比凤姐儿好搪塞些； 便添了一个探春，都想着不过是个未出闺阁的年轻小姐，
        且素日也最平和恬淡，因此都不在意，比凤姐儿前便懈怠了许多。 只三四天后，几 件事过手，渐觉探春精细处不让凤姐，只不过是言语安静、性情和顺而已。}
}
{
    \eA{Little, however, though one would think it, lady Feng was endowed with a
        poor physique. From her youth up, moreover, she had not known how to husband
        her health; and emulation and contentiousness had, more than anything else,
        combined to undermine her vital energies. Hence it was that although her
        complaint was a simple miscarriage, it had really, after all, been the
        outcome of loss of vigour.}
}
{
}

\Verse{6}
{
    \zA{可巧连 日有王公侯伯世袭官员十几处，皆系荣宁非亲即世交之家，或有升迁，或有黜降， 或有婚丧红白等事，王夫人贺吊迎送，应酬不暇，前边更无人照管。
        他二人便一日 皆在厅上起坐，宝钗便一日在上房监察，至王夫人回方散。 每于夜间针线暇时，临 寝之先，坐了轿，带领园中上夜人等，各处巡察一次。
        他三人如此一理，便觉比凤 姐儿当权时倒更谨慎了些。}
}
{
    \eA{After a month symptoms of emissions of blood began also to show themselves.
        And notwithstanding her reluctance to utter what she felt every one, at the
        sight of her sallow and emaciated face, readily concluded that she was not
        nursing herself as well as she should. Madame Wang therefore enjoined her
        merely to take her medicines and look to herself with due care;}
}
{
}

\Verse{7}
{
    \zA{因而里外下人都暗中抱怨，说：“刚刚的倒了一个‘巡 海夜叉’，又添了三个‘镇山太岁’，越发连夜里偷着吃酒玩的工夫都没了！”
        这日王夫人正是往锦乡侯府去赴席，李纨与探春早已梳洗，伺候出门去后，回 至厅上坐了。 刚吃茶时，只见吴新登的媳妇进来回说：“赵姨娘的兄弟赵国基昨儿
        出了事，已回过老太太、太太，说知道了，叫回姑娘来。” 说毕，便垂手旁侍，再 不言语。 彼时来回话者不少，都打听他二人办事如何。}
}
{
    \eA{and she would not allow her to disquiet her mind about the least thing. But
        (lady Feng) herself also gave way to misgivings lest her illness should
        assume some grave phase, and much though she laughed with one and all, she
        was ever mindful to steal time to attend to her health, feeling inwardly
        vexed at not being able to soon get back her old strength again.}
}
{
}

\Verse{8}
{
    \zA{若办得妥当，大家则安个畏 惧之心，若少有嫌隙不当之处，不但不畏服，一出二门，还说出许多笑话来取笑。
        吴新登的媳妇心中已有主意，若是凤姐前，他便早已献勤，说出许多主意、又查出 许多旧例来，任凤姐拣择施行； 如今他藐视李纨老实，探春是年轻的姑娘，所以只
        说出这一句话来，试他二人有何主见。 探春便问李纨，李纨想了一想，便道：“前 日袭人的妈死了，听见说赏银四十两，这也赏他四十两罢了。” 吴新登的媳妇听了，
        忙答应了个“是”，接了对牌就走。}
}
{
    \eA{But she had, as it happened, to dose herself with medicines and to nurse
        herself for three whole months, before she gradually began to rally and
        before the discharges stopped by degrees.}
}
{
}

\Verse{9}
{
    \zA{探春道：“你且回来。” 吴新登家的只得回来。 探春道：“你且别支银子。 我且问你：那几年老太太屋里的几位老姨奶奶，也有家 里的，也有外头的，有两个分别。
        家里的若死了人是赏多少？ 外头的死了人是赏多 少？ 你且说两个我们听听。” 一问，吴新登家的便都忘了，忙陪笑回说道：“这也
        不是什么大事，赏多赏少，谁还敢争不成？” 探春笑道：“这话胡闹。 依我说，赏 一百倒好!若不按理，别说你们笑话，明儿也难见你二奶奶。”}
}
{
    \eA{But we will abstain from any reference to these details which pertain to the
        future, suffice it now to add that though Madame Wang noticed her improved
        state, (she thought it) impossible for the time being for T'an Ch'un and Li
        Wan to resign their charge.}
}
{
}

\Verse{10}
{
    \zA{吴新登家的笑道： “既这么说，我查旧账去； 此时却不记得。” 探春笑道：“你办事办老了的，还不 记得，倒来难我们!你素日回你二奶奶，也现查去？
        若有这道理，凤姐姐还不算利害， 也就算是宽厚了。 还不快找了来我瞧!再迟一日，不说你们粗心，倒像我们没主意 了。” 吴新登家的满面通红，忙转身出来。
        众媳妇们都伸舌头。 这里又回别的事； 一时吴家的取了旧账来。 探春看时，两个家里的赏过皆二十 四两，两个外头的皆赏过四十两。}
}
{
    \eA{But so fidgetty was she lest with the large number of inmates in the garden
        proper control should not be exercised that she specially sent for Pao-ch'ai
        and begged of her to keep an eye over every place, explaining to her that
        the old matrons were of no earthly use, for whenever they could obtain any
        leisure, they drank and gambled;}
}
{
}

\Verse{11}
{
    \zA{外还有两个外头的，一个赏过一百两，一个赏过 六十两。 这两笔底下皆有原故：一个是隔省迁父母之柩，外赏六十两； 一个是现买 葬地，外赏二十两。
        探春便递给李纨看了，探春便说：“给他二十两银子，把这账 留下我们细看。” 吴新登家的去了。 忽见赵姨娘进来，李纨探春忙让坐。
        赵姨娘开口便说道：“这屋里的人，都踹 下我的头去还罢了，姑娘你也想一想，该替我出气才是！” 一面说，一面便眼泪鼻 涕哭起来。 探春忙道：“姨娘这话说谁？
        我竟不懂。}
}
{
    \eA{and slept during broad daylight, while they played at cards during the hours
        of night. "I know all about their doings," (she said). "When that girl Feng
        is well enough to go out, they have some little fear. But they're bound at
        present to consult again their own convenience. Yet you, dear child, are one
        in whom I can repose complete trust. Your brother and your female cousins
        are, on the one hand, young;}
}
{
}

\Verse{12}
{
    \zA{谁踹姨娘的头？ 说出来，我替姨 娘出气。” 赵姨娘道：“姑娘现踹我，我告诉谁去？” 探春听说，忙站起来说道： “我并不敢。” 李纨也忙站起来劝。
        赵姨娘道：“你们请坐下，听我说。 我这屋里 熬油似的熬了这么大年纪，又有你兄弟，这会子连袭人都不如了，我还有什么脸？ 连你也没脸面，别说是我呀。”
        探春笑道：“原来为这个，我说我并不敢犯法违礼。”}
}
{
    \eA{and I can, on the other, afford no spare time; so do exert yourself on my
        behalf for a couple of days, and exercise proper supervision. And should
        anything unexpected turn up, just come and tell it to me. Don't wait until
        our old lady inquires about it, as I shall then find myself in a corner with
        nothing to say in my defence. If those servants aren't on their good
        behaviour, mind you blow them up;}
}
{
}

\Verse{13}
{
    \zA{一面便坐了，拿账翻给赵姨娘瞧，又念给他听，又说道：“这是祖宗手里旧规矩， 人人都依着，偏我改了不成？ 这也不但袭人，将来环儿收了外头的，自然也是和袭
        人一样。 这原不是什么争大争小的事，讲不到有脸没脸的话上。 他是太太的奴才， 我是按着旧规矩办。 说办的好，领祖宗的恩典、太太的恩典；
        若说办的不公，那是 他糊涂不知福，也只好凭他抱怨去。 太太连房子赏了人，我有什么有脸的地方儿？ 一文不赏，我也没什么没脸的。}
}
{
    \eA{and if they don't listen to you, come and lay your complaint before me; for
        it will be best not to let anything assume a serious aspect." Pao-ch'ai
        listened to her appeal and felt under the necessity of volunteering to
        undertake the charge. The season was about the close of spring, so Tai-yü
        got her cough back again.}
}
{
}

\Verse{14}
{
    \zA{依我说，太太不在家，姨娘安静些，养神罢，何苦 只要操心？ 太太满心疼我，因姨娘每每生事，几次寒心。 我但凡是个男人，可以出
        得去，我早走了，立出一番事业来，那时自有一番道理，偏我是女孩儿家，一句多 话也没我乱说的。 太太满心里都知道，如今因看重我，才叫我管家务。 还没有做一
        件好事，姨娘倒先来作践我。 倘或太太知道了，怕我为难，不叫我管，那才正经没 脸呢!连姨娘真也没脸了！” 一面说，一面抽抽搭搭的哭起来。}
}
{
    \eA{But Hsiang-yün was likewise laid up in the Heng Wu Yüan, as she too was
        affected by the weather, and day after day she saw numberless doctors and
        took endless medicines. T'an Ch'un and Li Wan lived apart, but as they had
        of late assumed joint management of affairs, it was, unlike former years,
        extremely inconvenient even for the servants to go backwards and forwards to
        make their reports.}
}
{
}

\Verse{15}
{
    \zA{赵姨娘没话答对，便说道：“太太疼你，你该越发拉扯拉扯我们。 你只顾讨太 太的疼，就把我们忘了！” 探春道：“我怎么忘了？ 叫我怎么拉扯？
        这也问他们各人。 那一个主子不疼出力得用的人？ 那一个好人用人拉扯呢？” 李纨在旁只管劝说：“姨 娘别生气，也怨不得姑娘。
        他满心里要拉扯，口里怎么说的出来？” 探春忙道：“这 大嫂子也糊涂了!我拉扯谁？ 谁家姑娘们拉扯奴才了？ 他们的好歹，你们该知道，与 我什么相干？”}
}
{
    \eA{They consequently resolved that they should meet early every day in the
        small three-roomed reception-hall, at the south side of the garden gate, to
        transact what business there was, and that their morning meal over, they
        should after noon return again to their quarters.}
}
{
}

\Verse{16}
{
    \zA{赵姨娘气的问道：“谁叫你拉扯别人去了？ 你不当家，我也不来问 你。 你如今现在说一是一，说二是二!如今你舅舅死了，你多给了二三十两银子，
        难道太太就不依你？ 分明太太是好太太，都是你们尖酸克薄!可惜太太有恩无处使! 姑娘放心：这也使不着你的银子，明日等出了阁，我还想你额外照看赵家呢!如今
        没有长翎毛儿就忘了根本，只‘拣高枝儿飞’去了。” 探春没听完，气的脸白气噎， 越发呜呜咽咽的哭起来。 因问道：“谁是我舅舅？}
}
{
    \eA{This three-roomed hall had originally been got ready at the time of the
        visit of the imperial consort to her parents, to accommodate the attendants
        and eunuchs. This visit over, it proved, therefore, no longer of use, and
        the old matrons simply came to it every night to keep watch. But mild
        weather had now set in, and any complete fittings were quite superfluous.}
}
{
}

\Verse{17}
{
    \zA{我舅舅早升了九省的检点了!那里 又跑出一个舅舅来？ 我倒素昔按礼尊敬，怎么敬出这些亲戚来了!既这么说，每日环 儿出去，为什么赵国基又站起来？
        又跟他上学？ 为什么不拿出舅舅的款来？ 何苦来!谁 不知道我是姨娘养的，必要过两三个月寻出由头来，彻底来翻腾一阵，怕人不知道，
        故意表白表白!也不知道是谁给谁没脸!幸亏我还明白，但凡糊涂不知礼的，早急 了！” 李纨急得只管劝，赵姨娘只管还唠叨。 忽听有人说：“二奶奶打发平姑娘说话
        来了。”}
}
{
    \eA{All that could be seen about amounted to a few small pieces of furniture
        just sufficient for them to make themselves comfortable with. Over this hall
        was likewise affixed a placard, with the inscription in four characters:
        "Perfected philanthropy, published virtue!" Yet the place was generally
        known among the domestics as 'the discuss-matters-hall.'}
}
{
}

\Verse{18}
{
    \zA{赵姨娘听说，方把嘴止住。 只见平儿走来，赵姨娘忙陪笑让坐，又忙问： “你奶奶好些？ 我正要瞧去，就只没得空儿。” 李纨见平儿进来，因问他：“来作
        什么？” 平儿笑道：“奶奶说：赵姨奶奶的兄弟没了，恐怕奶奶和姑娘不知有旧例。 若照常例，只得二十两； 如今请姑娘裁度着，再添些也使得。”
        探春早已拭去泪痕， 忙说道：“又好好的添什么？ 谁又是二十四个月养的？ 不然，也是出兵放马、背着主 子逃出命来过的人不成？}
}
{
    \eA{To this hall, (Li Wan and T'an Ch'un) would daily adjourn at six in the
        morning, and leave it at noon, and the wives of the managers and other
        servants, who had any matters to lay before them, came and went in incessant
        strings. When the domestics heard that Li Wan would assume sole control,
        each and all felt secretly elated;}
}
{
}

\Verse{19}
{
    \zA{你主子真个倒巧，叫我开了例，他做好人，拿着太太不心 疼的钱，乐得做人情!你告诉他：我不敢添减混出主意。 他添他施恩，等他好了出 来，爱怎么添怎么添！”
        平儿一来时已明白了对半，今听这话越发会意。 见探春有 怒色，便不敢以往日喜乐之时相待，只一边垂手默侍。 时值宝钗也从上房中来，探春等忙起身让坐。
        未及开言，又有一个媳妇进来回 事，因探春才哭了，便有三四个小丫鬟捧了脸盆、巾帕、靶镜等物来。}
}
{
    \eA{for as Li Wan had always been considerate, forbearing and loth to inflict
        penalties, she would be, of course, they thought, easier to put off than
        lady Feng.}
}
{
}

\Verse{20}
{
    \zA{此时探春因 盘膝坐在矮板榻上，那捧盆丫鬟走至跟前，便双膝跪下，高捧脸盆，那两个丫鬟也 都在旁屈膝捧着巾帕并靶镜脂粉之饰。
        平儿见侍书不在这里，便忙上来与探春挽袖 卸镯，又接过一条大手巾来，将探春面前衣襟掩了，探春方伸手向脸盆中盥沐。 媳
        妇便回道：“奶奶，姑娘：家学里支环爷和兰哥儿一年的公费。” 平儿先道：“你 忙什么？ 你睁着眼看见姑娘洗脸，你不出去伺候着，倒先说话来!二奶奶跟前，你也
        这样没眼色来着？}
}
{
    \eA{Even when T'an Ch'un was added, they again remembered that she was only a
        youthful unmarried girl and that she too had ever shown herself goodnatured
        and kindly to a degree, so none of them worried their minds about her, and
        they became considerably more indolent than when they had to deal with lady
        Feng.}
}
{
}

\Verse{21}
{
    \zA{姑娘虽恩宽，我去回了二奶奶，只说你们眼里都没姑娘，你们都 吃了亏可别怨我。” 唬得那个媳妇忙陪笑说：“我粗心了！” 一面说，一面忙退出 去。
        探春一面匀脸，一面向平儿冷笑道：“你迟了一步，没见还有可笑的。 连吴姐 姐这么个办老了事的，也不查清楚了就来混我们。 幸亏我们问他，他竟有脸说‘忘
        了’，我说他回二奶奶事也忘了再找去？ 我料着你主子未必有耐性儿等他去找！” 平儿笑道：“他有这么一次，包管腿上的筋早折了两根。 姑娘别信他们。}
}
{
    \eA{But after the expiry of three or four days several concerns passed through
        her hands, which gave them an opportunity to gradually find out that T'an
        Ch'un did not, in smartness and thoroughness, yield to lady Feng, and that
        the only difference between them was that she was soft in speech and gentle
        in disposition.}
}
{
}

\Verse{22}
{
    \zA{那是他们 瞅着大奶奶是个菩萨，姑娘又是腼腆小姐，固然是托懒来混。” 说着，又向门外说 道：“你们只管撒野，等奶奶大安了，咱们再说。”
        门外的众媳妇都笑道：“姑娘， 你是个最明白的人，俗语说：‘一人作罪一人当。’ 我们并不敢欺蔽主子。 如今主 子是娇客，若认真惹恼了，死无葬身之地！”
        平儿冷笑道：“你们明白就好了。” 又陪笑向探春道：“姑娘知道，二奶奶本来事多，那里照看得这些？ 保不住不忽略。 俗语说：‘旁观者清。’}
}
{
    \eA{By a remarkable coincidence, princes, dukes, marquises, earls, and
        hereditary officials arrived for consecutive days from various parts; all of
        whom were, if not the relatives of the Jung and Ning mansions, at least
        their old friends. There were either those who had obtained transfers on
        promotion, or others who had been degraded;}
}
{
}

\Verse{23}
{
    \zA{这几年姑娘冷眼看着，或有该添该减的去处，二奶奶没行 到，姑娘竟一添减：头一件，与太太有益； 第二件，也不枉姑娘待我们奶奶的情义 了。”
        话未说完，宝钗李纨皆笑道：“好丫头，真怨不得凤丫头偏疼他!本来无可 添减之事，如今听你一说，倒要找出两件来斟酌斟酌，不辜负你这话。”
        探春笑道：“我一肚子气，正要拿他奶奶出气去，偏他碰了来，说了这些话， 叫我也没了主意了。”}
}
{
    \eA{either those, who had married, or those who had gone into mourning, and
        Madame Wang had so much congratulating and condoling, receiving and
        escorting to do that she had no time to attend to any entertaining. There
        was therefore less than ever any one in the front part to look after things.}
}
{
}

\Verse{24}
{
    \zA{一面说，一面叫进方才那媳妇来问：“环爷和兰哥家学里这 一年的银子，是做那一项用的？” 那媳妇便回说：“一年学里吃点心或者买纸笔，
        每位有八两银子的使用。” 探春道：“凡爷们的使用，都是各屋里月钱之内：环哥 的是姨娘领二两； 宝玉的，老太太屋里袭人领二两； 兰哥儿是大奶奶屋里领：怎么
        学里每人多这八两？ 原来上学去的是为这八两银子!从今日起，把这一项蠲了。 平儿 回去，告诉你奶奶，说我的话，把这一条务必免了。” 平儿笑道：“早就该免。}
}
{
    \eA{So while (T'an Ch'un and Li Wan) spent their whole days in the hall,
        Pao-ch'ai tarried all day in the drawing-rooms, to keep an eye over what was
        going on; and they only betook themselves back to their quarters after
        Madame Wang's return. Of a night, they whiled away their leisure hours by
        doing needlework;}
}
{
}

\Verse{25}
{
    \zA{旧 年奶奶原说要免来着，因年下忙，就忘了。” 那媳妇只得答应着去了。 就有大观园中媳妇捧了饭盒子来，侍书素云早已抬过一张小饭桌来，平儿也忙 着上菜。
        探春笑道：“你说完了话，干你的去罢，在这里又忙什么？” 平儿笑道： “我原没事，二奶奶打发了我来，一则说话，二则怕这里的人不方便，叫我帮着妹
        妹们伏侍奶奶姑娘来了。” 探春因问：“宝姑娘的怎么不端来一处吃？”}
}
{
    \eA{but they would, previous to retiring to sleep, get into their chairs, and,
        taking along with them the servants, whose duty it was to be on night watch
        in the garden, and other domestics as well, they visited each place on their
        round. Such was the control exercised by these three inmates that signs were
        not wanting to prove that greater severity was observed than in the days
        when the management devolved on lady Feng.}
}
{
}

\Verse{26}
{
    \zA{丫鬟们听 说，忙出至檐外，命媳妇们去说：“宝姑娘如今在厅上一处吃，叫他们把饭送了这 里来。”
        探春听说，便高声说道：“你别混支使人!那都是办大事的管家娘子们， 你们支使他要饭要茶的？ 连个高低都不知道!平儿这里站着，叫他叫去。” 平儿忙答
        应了一声出来，那些媳妇们都悄悄的拉住笑道：“那里用姑娘去叫？ 我们已有人叫 去了。” 一面说，一面用绢子掸台阶的土，说：“姑娘站了半天，乏了，这太阳地
        里歇歇儿罢。” 平儿便坐下。}
}
{
    \eA{To this reason must be assigned the fact that all the servants attached
        inside as well as outside cherished a secret grudge against them. "No
        sooner," they insinuated, "has one patrolling ogre come than they add three
        more cerberean sort of spring josses so that even at night we've got less
        time than ever to sip a cup of wine and indulge in a romp!"}
}
{
}

\Verse{27}
{
    \zA{又有茶房里的两个婆子拿了个坐褥铺下，说：“石头 冷，这是极干净的，姑娘将就坐一坐儿罢。” 平儿点头笑道：“多谢。” 一个又捧
        了一碗精致新茶出来，也悄悄笑说：“这不是我们常用的茶，原是伺候姑娘们的， 姑娘且润一润罢。” 平儿遂欠身接了，因指众媳妇悄悄说道：“你们太闹的不像了。
        他是个姑娘家，不肯发威动怒，这是他尊重，你们就藐视欺负他。}
}
{
    \eA{On the day that Madame Wang was going to a banquet at the mansion of the
        Marquis of Chin Hsiang, Li Wan and T'an Ch'un arranged their coiffure and
        performed their ablutions at an early hour; and after waiting upon her until
        she went out of doors, they repaired into the hall and installed themselves
        in their seats. But just as they were sipping their tea, they espied Wu
        Hsin-teng's wife walk in. "Mrs.}
}
{
}

\Verse{28}
{
    \zA{果然招他动了大 气，不过说他一个粗糙就完了，你们就现吃不了的亏!他撒个娇儿，太太也得让他 一二分，二奶奶也不敢怎么。
        你们就这么大胆子小看他，可是鸡蛋往石头上碰。” 众人都忙道：“我们何尝敢大胆了？ 都是赵姨娘闹的。” 平儿也悄悄的道：“罢了!
        好奶奶们，‘墙倒众人推’。 那赵姨娘原有些颠倒，着三不着两，有了事就都赖他。 你们素日那眼里没人、心术利害，我这几年难道还不知道!二奶奶要是略差一点儿
        的，早叫你们这些奶奶们治倒了。}
}
{
    \eA{Chao's brother, Chao Kuo-chi," she observed, "departed this life yesterday;
        the tidings have already been reported to our old mistress and our lady, who
        said that it was all right, and bade me tell you, Miss." At the close of
        this announcement, she respectfully dropped her arms against her body, and
        stood aloof without adding another word.}
}
{
}

\Verse{29}
{
    \zA{饶这么着，得一点空儿，还要难他一难，好几次 没落了你们的口声。 众人都说他利害，你们都怕他，惟我知道他心里也就不算不怕 你们的。
        前儿我们还议论到这里：再不能依头顺尾，必有两场气生。 那三姑娘虽是 个姑娘，你们都横看了他!二奶奶在这些大姑子小姑子里头，也就只单怕他五分儿。
        你们这会子倒不把他放在眼里了。” 正说着，只见秋纹走来，众媳妇忙赶着问好，又说：“姑娘也且歇歇，里头摆 饭呢。 等撤下桌子来，再回话去罢。”}
}
{
    \eA{The servants, who came at this season to lay their reports before (T'an
        Ch'un and Li Wan), mustered no small number. But they all endeavoured to
        find out how their two new mistresses ran the household;}
}
{
}

\Verse{30}
{
    \zA{秋纹笑道：“我比不得你们，我那里等得？” 说着，便直要上厅去。 平儿忙叫：“快回来！” 秋纹回头，见了平儿，笑道：“你
        又在这里充什么‘外围子的防护’？” 一面回身便坐在平儿褥上。 平儿悄问：“回 什么？” 秋纹道：“问一问宝玉的月钱、我们的月钱，多早晚才领？”
        平儿道：“这 什么大事!你快回去告诉袭人，说我的话：凭有什么事，今日都别回。 若回一件管 驳一件，回一百件管驳一百件。”
        秋纹听了，忙问：“这是为什么？”}
}
{
    \eA{for as long they managed things properly, one and all willingly resolved to
        respect them, but in the event of the least disagreement or improper step,
        not only did they not submit to them, but they also spread, the moment they
        put their foot outside the second gate, numberless jokes on their account
        and made fun of them. Wu Hsin-teng's wife had thus devised an experiment in
        her own mind.}
}
{
}

\Verse{31}
{
    \zA{平儿与众媳 妇等都忙告诉他原故，又说：“正要找几处利害事与有体面的人来开例，作法子镇 压，与众人作榜样呢。 何苦你们先来碰在这钉子上？
        你这一去说了，他们若拿你们 也作一二件榜样，又碍着老太太、太太； 若不拿着你们做一二件，人家又说：‘偏
        一个向一个，仗着老太太、太太威势的就怕，不敢惹，只拿着软的做鼻子头。’ 你 听听罢，二奶奶的事他还要驳两件，才压得众人口声呢。”}
}
{
    \eA{Had she had to deal with lady Feng, she would have long ago made an attempt
        to show off her zeal by proposing numerous alternatives and discovering
        various bygone precedents, and then allowed lady Feng to make her own choice
        and take action;}
}
{
}

\Verse{32}
{
    \zA{秋纹听了，伸了伸舌头 笑道：“幸而平姐姐在这里，没得臊一鼻子灰，趁早知会他们去。” 说着便起身走 了。 接着宝钗的饭至，平儿忙进来伏侍。
        那时赵姨娘已去，三人在板床上吃饭，宝 钗面南，探春面西，李纨面东。 众媳妇皆在廊下静候，里头只有他们紧跟常侍的丫 鬟伺候，别人一概不敢擅入。
        这些媳妇们都悄悄的议论说：“大家省事罢，别安着 没良心的主意。 连吴大娘才都讨了没意思，咱们又是什么有脸的？” 都一边悄议， 等饭完回事。}
}
{
    \eA{but, in this instance, she looked with such disdain on Li Wan, on account of
        her simplicity, and on T'an Ch'un, on account of her youthfulness, that she
        volunteered only a single sentence, in order to put both these ladies to the
        test, and see what course they would be likely to adopt. "What shall we do?"
        T'an Ch'un asked Li Wan. Li Wan reflected for a while.}
}
{
}

\Verse{33}
{
    \zA{此时里面惟闻微嗽之声，不闻碗箸之响。 一时，只见一个丫头将帘栊 高揭，又有两个将桌抬出。 茶房内有三个丫鬟，捧着三个沐盆儿，见饭桌已出，三 人便进去了。
        一回又捧出沐盆并漱盂来，方有侍书、素云、莺儿三个人，每人用茶 盘捧了三盖碗茶进去。 一时等他三人出来，侍书命小丫头子：“好生伺候着，我们
        吃饭来换你们，可又别偷坐着去。” 众媳妇们方慢慢的安分回事，不敢如先前轻慢 疏忽了。}
}
{
    \eA{"The other day," she rejoined, "that Hsi Jen's mother died, I heard that she
        was given forty taels. So now give her forty taels as well and have done!"
        Upon hearing this proposal, Wu Hsin-teng's wife eagerly expressed her
        acquiescence, by uttering a yes; and taking over the permit she was going on
        her way at once. "Come back," shouted T'an Ch'un. "Wu Hsing-teng's wife had
        perforce to retrace her footsteps.}
}
{
}

\Verse{34}
{
    \zA{探春气方渐平，因向平儿道：“我有一件大事，早要和你奶奶商议，如今可巧 想起来。 你吃了饭快来。 宝姑娘也在这里，咱们四个人商议了，再细细的问你奶奶
        可行可止。” 平儿答应回去。 凤姐因问：“为何去这半日？” 平儿便笑着将方才的 原故细细说与他听了。
        凤姐儿笑道：“好，好，好!好个三姑娘，我说不错，只可 惜他命薄，没托生在太太肚里。” 平儿笑道：“奶奶也说糊涂话了。 他就不是太太
        养的，难道谁敢小看他，不和别的一样看待么？”}
}
{
    \eA{"Wait, don't get the money yet," T'an Ch'un remarked. "I want to ask you
        something. Some of the old secondary wives, attached years back to our
        venerable senior's rooms, lived inside the establishment; others outside;
        there were these two distinctions between them. Now if any of them died at
        home, how much was allowed them? And how much was allotted to such as died
        outside?}
}
{
}

\Verse{35}
{
    \zA{凤姐叹道：“你那里知道？ 虽然 正出庶出是一样，但只女孩儿却比不得儿子。 将来作亲时，如今有一种轻狂人，先 要打听姑娘是正出是庶出，多有为庶出不要的。
        殊不知庶出只要人好，比正出的强 百倍呢。 将来不知那个没造化的，为挑正庶误了事呢，也不知那个有造化的，不挑 正庶的得了去。”
        说着，又向平儿笑道：“你知道我这几年生了多少省俭的法子， 一家子大约也没个背地里不恨我的。 我如今也是骑上老虎了，虽然看破些，无奈一 时也难宽放。}
}
{
    \eA{Tell us what was given in either case for our guidance." As soon as Wu
        Hsin-teng's wife was asked this question, every detail bearing on the
        subject slipped from her memory. Hastily forcing a smile, "This is," she
        replied, "nothing of any such great consequence. Whether much or little be
        allowed, who'll ever venture to raise a quarrel about it?" T'an Ch'un then
        smiled. "This is all stuff and nonsense!" she exclaimed.}
}
{
}

\Verse{36}
{
    \zA{二则家里出去的多，进来的少，凡有大小事儿，仍是照着老祖宗手里 的规矩。 却一年进的产业又不及先时多，省俭了外人又笑话，老太太、太太也受委
        屈，家下也抱怨克薄。 若不趁早儿料理省俭之计，再几年就都赔尽了。” 平儿道： “可不是这话!将来还有三四位姑娘，还有两三个小爷们，一位老太太，这几件大
        事未完呢。” 凤姐儿笑道：“我也虑到这里，倒也够了，宝玉和林妹妹，他两个一 娶一嫁，可以使不着官中钱，老太太自有体己拿出来。}
}
{
    \eA{"My idea is that it would be better to give a hundred taels. For if we don't
        comply with what's right, we shall, not to speak of your ridiculing us, find
        it also a hard job by and bye to face your mistress Secunda." "Well, in that
        case," laughed Wu Hsin-teng's wife, "I'll go and look up the old accounts. I
        can't recollect anything about them just at this moment."}
}
{
}

\Verse{37}
{
    \zA{二姑娘是大老爷那边的，也 不算。 剩了三四个，满破着每人花上七八千银子。 环哥娶亲有限，花上三千银子， 若不够，那里省一抿子也就够了。
        老太太的事出来，一应都是全了的，不过零星杂 项使费些，满破三五千两。 如今再俭省些，陆续就够了。 只怕如今平空再生出一两 件事来，可就了不得了。
        咱们且别虑后事，你且吃了饭，快听他们商议什么。 这正 碰了我的机会，我正愁没个膀臂。 虽有个宝玉，他又不是这里头的货，纵收伏了他 也不中用。}
}
{
    \eA{"You're quite an old hand in the management of affairs," T'an Ch'un observed
        with a significant smile, "and can't you remember, but come instead to
        perplex us? Whenever you've had anything of the kind to lay before your lady
        Secunda, have you also had to go first and look it up? But if this has been
        the practice, lady Feng can't be looked upon as being such a dreadful
        creature.}
}
{
}

\Verse{38}
{
    \zA{大奶奶是个佛爷，也不中用。 二姑娘更不中用，亦且不是这屋里的人。 四姑娘小呢。 兰小子和环儿更是个燎毛的小冻猫子，只等有热灶火炕让他钻去罢，
        ――真真一个娘肚子里跑出这样天悬地隔的两个人来，我想到那里就不服!再者林 丫头和宝姑娘他两个人倒好，偏又都是亲戚，又不好管咱们家务事。 况且一个是美
        人灯儿，风吹吹就坏了； 一个是拿定了主意，不干己事不张口，一问摇头三不知， 也难十分去问他。}
}
{
    \eA{One could very well call her lenient and kind. Yet don't you yet hurry to go
        and hunt them up and bring them to me to see? If we dilly-dally another day,
        they won't run you people down for your coarse-mindedness, but we will seem
        to have been driven to our wits' ends!" Wu Hsin-teng's wife got quite
        scarlet in the face. Promptly twisting herself round, she quitted the hall;}
}
{
}

\Verse{39}
{
    \zA{倒只剩了三姑娘一个，心里嘴里都也来得，又是咱家的正人，太 太又疼他，虽然脸上淡淡的，皆因是赵姨娘那老东西闹的，心里却是和宝玉一样呢。
        比不得环儿，实在令人难疼，要依我的性子，早撵出去了!如今他既有这主意，正 该和他协同，大家做个膀臂，我也不孤不独了。 按正礼天理良心上论，咱们有他这
        一个人帮着，咱们也省些心，与太太的事也有益。 若按私心藏奸上论，我也太行毒 了，也该抽回退步，回头看看；}
}
{
    \eA{while the whole bevy of married women stretched out their tongues
        significantly. During her absence, other matters were reported. But in a
        little while, Wu Hsin-teng's wife returned with the old accounts.}
}
{
}

\Verse{40}
{
    \zA{再要穷追苦克，人恨极了，他们笑里藏刀，咱们两 个才四个眼睛两个心，一时不防，倒弄坏了。 趁着紧溜之中，他出头一料理，众人 就把往日咱们的恨暂可解了。
        还有一件，我虽知你极明白，恐怕你心里挽不过来， 如今嘱咐你：他虽是姑娘家，心里却事事明白，不过是言语谨慎。 他又比我知书识 字，更利害一层了。
        如今俗语说：‘擒贼必先擒王。’ 他如今要作法开端，一定是 先拿我开端，倘或他要驳我的事，你可别分辩，你只越恭敬越说驳的是才好。}
}
{
    \eA{On inspection, T'an Ch'un found that for a couple of secondary wives, who
        had lived in the establishment, twenty-four taels had been granted, and that
        for two, whose quarters had been outside, forty taels had in each case been
        allowed. Besides these two, others were mentioned, who had lived outside the
        mansion; to one of whom a hundred taels had been given, and to the other,
        sixty taels.}
}
{
}

\Verse{41}
{
    \zA{千万 别想着怕我没脸，和他一强，就不好了。” 平儿不等说完，便笑道：“你太把人看糊涂了!我才已经行在先了，这会子才 嘱咐我。”
        凤姐儿笑道：“我是恐怕你心里眼里只有了我、一概没有他人之故，不 得不嘱咐，既已行在先，更比我明白了。 这不是你又急了，满嘴里‘你’呀‘我’ 的起来了！”
        平儿道：“偏说‘你’!你不依，这不是嘴巴子？ 再打一顿。 难道这脸 上还没尝过的不成？” 凤姐儿笑道：“你这小蹄子儿，要掂多少过儿才罢？}
}
{
    \eA{Under these two records, the reasons were assigned. In the one case, the
        coffins of father and mother had had to be removed from another province,
        and sixty taels extra had consequently been granted. In the other, an
        additional twenty taels had been allowed, as a burial-place had to be
        purchased at the time. T'an Ch'un handed the accounts to Li Wan for her
        perusal.}
}
{
}

\Verse{42}
{
    \zA{你看我 病的这个样儿，还来怄我呢。 过来坐下，横竖没人来，咱们一处吃饭是正经。” 说 着，丰儿等三四个小丫头子进来，放小炕桌。
        凤姐只吃燕窝粥，两碟子精致小菜， 每日分例菜已暂减去。 丰儿便将平儿的四样分例菜端至桌上，与平儿盛了饭来。 平
        儿屈一膝于炕沿之上，半身犹立于炕下，陪着凤姐儿吃了饭，伏侍漱口毕，吩咐了 丰儿些话，方往探春处来。 只见院中寂静，人已散出。
        要知后事何如，且听下回分解。 (本章完)}
}
{
    \eA{"Give her twenty taels," readily suggested T'an Ch'un. "Leave these accounts
        here for us to examine minutely." Wu Hsin-teng's wife then walked away. But
        unexpectedly Mrs. Chao entered the hall. Li Wan and T'an Ch'un speedily
        pressed her to take a seat. Mrs. Chao then broke the silence. "All the
        inmates of these rooms have trampled me under heel," she said, "but never
        mind!}
}
{
}

\Verse{43}
{
    \zA{}
}
{
    \eA{Yet, my child, just ponder, it is only fair that you should take my part."
        While ventilating her grievances, her eyes got moist, her nose watered, and
        she began to sob. "To whom are you alluding Mrs. Chao?" T'an Ch'un hastily
        inquired. "I can't really make out what you're driving at. Who tramples you
        under foot? Speak out and I'll take up your cudgels." "You're now trampling
        me down yourself, young lady," Mrs. Chao observed. "And to whom can I go and
        tell my grievance?" T'an Ch'un, at these words, jumped up with alacrity. "I
        never would presume to do any such thing," she protested. Li Wan too
        vehemently sprung to her feet to proffer her some good counsel. "Pray seat
        yourselves, both of you," Mrs. Chao cried, "and listen to what I have to
        say.}
}
{
}

\Verse{44}
{
    \zA{}
}
{
    \eA{I've had, like simmering oil, to consume away in these rooms to this
        advanced age. There's also your brother besides. Yet I can't compare myself
        now even to Hsi Jen, and what credit do I enjoy? But you haven't as well any
        face, so don't let's speak of myself." "It was really on account of this,"
        T'an Ch'un smiled, "that I said that I didn't presume to disregard right and
        to violate propriety." While she spoke, she resumed her seat, and taking up
        the accounts, she turned them over for Mrs. Chao to glance at, after which
        she read them out to her for her edification. "These are old customs," she
        proceeded, "enforced by the seniors of the family, and every one complies
        with them, and could I ever, pray, have changed them? These will hold good
        not only with Hsi Jen;}
}
{
}

\Verse{45}
{
    \zA{}
}
{
    \eA{but even when by and bye Huan-erh takes a concubine, the same course will
        naturally be adopted as in the case of Hsi Jen. This is no question for any
        large quarrels or small disputes, and no mention should be made about face
        or no face. She's our Madame Wang's servant-girl, and I've dealt with her
        according to a long-standing precedent. Those who say that I've taken
        suitable action will come in for our ancestors' bounty and our lady's bounty
        as well. But should any one uphold that I've adopted an unfair course, that
        person is devoid of all common sense and totally ignorant of what a blessing
        means. The only thing she can do is to foster as much resentment as she
        chooses. Our lady, Madame Wang, may even give a present of a house to any
        one;}
}
{
}

\Verse{46}
{
    \zA{}
}
{
    \eA{what credit is that to me? Again, she may not give a single cash, but even
        that won't imply any loss of face, as far as I am concerned. What I have to
        say is that as Madame Wang is away from home, you should quietly look after
        yourself a bit. What's the good of worrying and fretting? Our lady is
        extremely fond of me; and, if, at different times, a chilliness has sprung
        up on her part, it's because you, Mrs. Chao, have again and again been
        officious. Had I been a man and able to have gone abroad, I would long ago
        have run away and started some business. I would then have had something of
        my own to attend to. But, as it happens, I am a girl, so that I can't even
        recklessly utter so much as a single remark.}
}
{
}

\Verse{47}
{
    \zA{}
}
{
    \eA{Madame Wang is well aware of it in her heart. And it's now because she
        entertains a high opinion of me that she recently bade me assume the charge
        of domestic affairs. But before I've had time enough to do a single good
        act, here you come, Mrs. Chao, to lay down the law. If this reaches Madame
        Wang's ear, I fear I shall get into trouble. She won't let me exercise any
        control, and then I shall, in real earnest, come in for no face. But even
        you, Mrs. Chao, will then actually lose countenance." Reasoning with her,
        she so little could repress her tears that they rolled down her cheeks. Mrs.
        Chao had not a word more to say to refute her arguments with.}
}
{
}

\Verse{48}
{
    \zA{}
}
{
    \eA{"If Madame Wang loves you," she simply responded, "there's still more reason
        why you should have drawn us into her favour. (Instead of that), all you
        think about is to try and win Madame Wang's affections, and you forget all
        about us." "How ever did I forget you?" T'an Ch'un exclaimed. "How would you
        have me drag you into favour? Go and ask every one of them, and you'll see
        what mistress is indifferent to any one, who exerts her energies and makes
        herself useful, and what worthy person requires being drawn into favour?" Li
        Wan, who stood by, did her best to pacify them with her advice. "Mrs. Chao,"
        she argued, "don't lose your temper! Neither should you feel any ill-will
        against this young lady of yours.}
}
{
}

\Verse{49}
{
    \zA{}
}
{
    \eA{Had she even at heart every good intention to lend you a hand, how could she
        put it into words?" "This worthy senior dame," T'an Ch'un impatiently
        interposed, "has also grown quite dense! Whom could I drag into favour? Why,
        in what family, do the young ladies give a lift to slave-girls? Their
        qualities as well as defects should all alike be well known to you people.
        And what have they got to do with me?" Mrs. Chao was much incensed. "Who
        tells you," she asked, "to give a lift to any one? Were it not that you
        looked after the house, I wouldn't have come to inquire anything of you. But
        anything you may suggest is right; so had you, now that your maternal uncle
        is dead, granted twenty or thirty taels in excess, is it likely that Madame
        Wang would not have given you her consent?}
}
{
}

\Verse{50}
{
    \zA{}
}
{
    \eA{It's evident that our Madame Wang is a good woman and that it's you people
        who are mean and stingy. Unfortunately, however, her ladyship has with all
        her bounty no opportunity of exercising it. You could, my dear girl, well
        set your mind at ease. You wouldn't, in this instance, have had to spend any
        of your own money; and at your marriage by and bye, I would still have borne
        in mind the exceptional regard you had shown the Chao family. But now that
        you've got your full plumage, you've forgotten your extraction, and chosen a
        lofty branch to fly to." Before T'an Ch'un had heard her to the end, she
        flew into such a rage that her face blanched; and choking for breath, she
        gasped and panted.}
}
{
}

\Verse{51}
{
    \zA{}
}
{
    \eA{Sobbing, she asked the while: "Who's my maternal uncle? My maternal uncle
        was at the end of the year promoted to be High Commissioner of the Nine
        Provinces! How can another maternal uncle have cropped up? It's because I've
        ever shown that reverence enjoined by the rites that other relatives have
        now more than ever turned up. If what you say be the case, how is it that
        every day that Huan-erh goes out, Chao Kuo-chi too stands up, and follows
        him to school? Why doesn't he put on the airs of an uncle? What's the reason
        that he doesn't? Who isn't aware of the fact that I'm born of a concubine?
        Would it require two or three months' time to trace my extraction?}
}
{
}

\Verse{52}
{
    \zA{}
}
{
    \eA{But the fact is you've come to kick up all this hullaballoo for fear lest
        people shouldn't be alive to the truth; and with the express design of
        making it public all over the place! But I wonder who of us two will make
        the other lose face? Luckily, I've got my wits about me; for had I been a
        stupid creature ignorant of good manners, I would long ago have lost all
        patience." Li Wan was much concerned, but she had to continue to exhort them
        to desist. But Mrs. Chao proceeded with a long rigmarole until a servant was
        unexpectedly heard to report that lady Secunda had sent Miss Ping to deliver
        a message. Mrs. Chao caught the announcement, and eventually held her peace,
        when they espied P'ing erh making her appearance.}
}
{
}

\Verse{53}
{
    \zA{}
}
{
    \eA{Mrs. Chao hastily forced a saturnine smile, and motioned to her to take a
        seat. "Is your lady any better?" she went on to inquire with vehemence. "I
        was just thinking of going to look her up; but I could find no leisure!"
        Upon seeing P'ing Erh enter, Li Wan felt prompted to ask her the object of
        her visit. "My lady says," P'ing Erh smilingly responded, "that she
        apprehends, now that Mrs. Chao's brother is dead, that your ladyship and
        you, miss, are not aware of the existence of an old precedent. According to
        the ordinary practice no more need be given than twenty taels; but she now
        requests you, miss, to consider what would be best to do; if even you add a
        good deal more, it will do well enough."}
}
{
}

\Verse{54}
{
    \zA{}
}
{
    \eA{T'an Ch'un at once wiped away all traces of tears. "What's the use of
        another addition, when there's no valid reason for it?" she promptly
        demurred. "Who has again been twenty months in the womb? Or is it forsooth
        any one who's gone to the wars, and managed to escape with his life,
        carrying his master on his back? Your mistress is certainly very ingenious!
        She tells me to disregard the precedent, in order that she should pose as a
        benefactress! She wishes to take the money, which Madame Wang spurns, so as
        to reap the pleasure of conferring favours! Just you tell her that I could
        not presume to add or reduce anything, or even to adopt any reckless
        decision. Let her add what she wants and make a display of bounty.}
}
{
}

\Verse{55}
{
    \zA{}
}
{
    \eA{When she gets better and is able to come out, she can effect whatever
        additions she fancies." The moment P'ing Erh arrived, she obtained a fair
        insight (into lady Feng's designs), so when she heard the present remarks,
        she grasped a still more correct idea of things. But perceiving an angry
        look about T'an Ch'un's face, she did not have the temerity to behave
        towards her as she would, had she found her in the high spirits of past
        days. All she did therefore was to stand aloof with her arms against her
        sides and to wait in rigid silence. Just at that moment, however, Pao-ch'ai
        dropped in, on her return from the upper rooms.}
}
{
}

\Verse{56}
{
    \zA{}
}
{
    \eA{T'an Ch'un quickly rose to her feet, and offered her a seat. But before they
        had had time to exchange any words, a married woman likewise came to report
        some business. But as T'an Ch'un had been having a good cry, three or four
        young maids brought her a basin, towel, and hand-glass and other articles of
        toilette. T'an Ch'un was at the moment seated cross-legged, on a low wooden
        couch, so the maid with the basin had, when she drew near, to drop on both
        her knees and lift it high enough to bring it within reach. The other two
        girls prostrated themselves next to her and handed the towels and the rest
        of the toilet things, which consisted of a looking-glass, rouge and powder.}
}
{
}

\Verse{57}
{
    \zA{}
}
{
    \eA{But P'ing Erh noticed that Shih Shu was not in the room, and approaching
        T'an Ch'un with hasty step, she tucked up her sleeves for her and unclasped
        her bracelets. Seizing also a large towel from the hands of one of the
        maids, she covered the lapel on the front part of T'an Ch'un's dress;
        whereupon T'an Ch'un put out her hands, and washed herself in the basin. "My
        lady and miss," the married woman observed, "may it please you to pay what
        has been spent in the family school for Mr. Chia Huan and Mr.. Chia Lan
        during the year." P'ing Erh was the first to speak. "What are you in such a
        hurry for?" she cried. "You've got your eyes wide open, and must be able to
        see our young lady washing her face; instead of coming forward to wait on
        her, you start talking!}
}
{
}

\Verse{58}
{
    \zA{}
}
{
    \eA{Do you also behave in this blind sort of way in the presence of your lady
        Secunda? This young lady is, it's true, generous and lenient, but I'll go
        and report you to your mistress. I'll simply tell her that you people have
        no eye for Miss T'an Ch'un. But when you find yourselves in a mess, don't
        bear me any malice." At this hint the woman took alarm, and hastily forcing
        a smile, she pleaded guilty. "I've been rude," she exclaimed. With these
        words, she rushed with all despatch out of the room. T'an Ch'un smoothed her
        face. While doing so, she turned herself towards P'ing Erh and gave her a
        cynical smile. "You've come just one step too late," she remarked. "You
        weren't in time to see something laughable!}
}
{
}

\Verse{59}
{
    \zA{}
}
{
    \eA{Even sister Wu, an old hand at business though she be, failed to look up
        clearly an old custom and came to play her tricks on us. But when we plied
        her with questions, she luckily had the face to admit that it had slipped
        from her memory. 'Do you,' I insinuated, 'also forget, when you've got
        anything to report to lady Secunda? and have you subsequently to go and hunt
        up all about it?' Your mistress can't, I fancy, be so patient as to wait
        while she goes and institutes proper search." P'ing Erh laughed. "Were she
        to have behaved but once in this wise," she observed, "I feel positive that
        a couple of the tendons of her legs would have long ago been snapped. But,
        Miss, don't credit all they say.}
}
{
}

\Verse{60}
{
    \zA{}
}
{
    \eA{It's because they see that our senior mistress is as sweet-tempered as a
        'P'u-sa,' and that you, miss, are a modest young lady, that they, naturally,
        shirk their duties and come and take liberties with you. Your mind is set
        upon playing the giddy dogs," continuing, she added; speaking towards those
        beyond the doorway; "but when your mistress gets quite well again, we'll
        tell her all." "You're gifted with the greatest perspicacity, miss," the
        married women, standing outside the door, smiled in chorus. "The proverb
        says: 'the person who commits a fault must be the one to suffer.' We don't
        in any way presume to treat any mistress with disdain. Our mistress at
        present is in delicate health, and if we intentionally provoke her, may we,
        when we die, have no place to have our corpses interred in."}
}
{
}

\Verse{61}
{
    \zA{}
}
{
    \eA{P'ing Erh laughed a laugh full irony. "So long as you're aware of this, it's
        well and good," she said. And smiling a saturnine smile, she resumed,
        addressing herself to T'an Ch'un: "Miss, you know very well how busy our
        lady has been and how little she could afford the time to keep this tribe of
        people in order. Of course, they couldn't therefore, be prevented from
        becoming remiss. The adage has it: 'Lookers-on are clear of sight!'}
}
{
}

\Verse{62}
{
    \zA{}
}
{
    \eA{During all these years that you, have looked on dispassionately, there have
        possibly been instances on which, though additions or reductions should have
        been made, our lady Secunda has not been able to effect them, so, miss, do
        add or curtail whatever you may deem necessary, in order that, first, Madame
        Wang may be benefited, and that, secondly, you mayn't too render nugatory
        the kindness with which you ever deal towards our mistress." But scarcely
        had she finished, than Pao-ch'ai and Li Wan smilingly interposed. "What a
        dear girl!" they ejaculated. "One really can't feel angry with that hussy
        Feng for being partial to her and fond of her.}
}
{
}

\Verse{63}
{
    \zA{}
}
{
    \eA{We didn't, at first, see how we could very well alter anything by any
        increase or reduction, but after what you've told us, we must hit upon one
        or two things and try and devise means to do something, with a view of not
        showing ourselves ungrateful of the advice you've tendered us." "My heart
        was swelling with indignation," T'an Ch'un observed laughing, "and I was
        about to go and give vent to my temper with her mistress, but now that she
        (P'ing Erh) has happened to come, she has, with a few words, quite dissuaded
        me from my purpose." While she spoke, she called the woman, who had been
        with them a few minutes back, to return into the room.}
}
{
}

\Verse{64}
{
    \zA{}
}
{
    \eA{"For what things for Mr. Chia Huan and Mr. Chia Lau was the money expended
        during the year in the family school?" she inquired of her. "For cakes,"
        replied the woman, "they ate during the year at school; or for the purchase
        of paper and pens. Each one of them is allowed eight taels." "The various
        expenses on behalf of the young men," T'an Ch'un added, "are invariably paid
        in monthly instalments to the respective households. For cousin Chia Huan's,
        Mrs. Chao receives two taels. For Pao-yü's, Hsi Jen draws two taels from our
        venerable senior's suite of apartments. For cousin Chia Lan's, some one, in
        our senior lady's rooms, gets the proper allowance.}
}
{
}

\Verse{65}
{
    \zA{}
}
{
    \eA{So how is it that these extra eight taels have to be disbursed at school for
        each of these young fellows? Is it really for these eight taels that they go
        to school? But from this day forth I shall put a stop to this outlay. So
        P'ing Erh, when you get back, tell your mistress that I say that this item
        must absolutely be done away with." "This should have been done away with
        long ago," P'ing Erh smiled. "Last year our lady expressed her intention to
        eliminate it, but with the endless things that claimed her attention about
        the fall of the year, she forgot all about it." The woman had no other
        course than to concur with her views and to walk away. But the married women
        thereupon arrived from the garden of Broad Vista with the boxes of eatables.}
}
{
}

\Verse{66}
{
    \zA{}
}
{
    \eA{So Shih Shu and Su Yün at once brought a small dining-table, and P'ing Erh
        began to fuss about laying the viands on it. "If you've said all you had,"
        T'an Ch'un laughed, "you'd better be off and attend to your business. What's
        the use of your bustling about here?" "I've really got nothing to do," P'ing
        Erh answered smiling. "Our lady Secunda sent me first, to deliver a message;
        and next, because she feared that the servants in here weren't handy enough.
        The fact is, she bade me come and help the girls wait on you, my lady, and
        on you, miss." "Why don't you bring Mrs. Pao's meal so that she should have
        it along with us?" T'an Ch'un then inquired. As soon as the waiting-maids
        heard her inquiry, they speedily rushed out and went under the eaves.}
}
{
}

\Verse{67}
{
    \zA{}
}
{
    \eA{"Go," they cried, directing the married women, "and say that Miss Pao-ch'ai
        would like to have her repast just now in the hall along with the others,
        and tell them to send the eatables here." T'an Ch'un caught their
        directions. "Don't be deputing people to go on reckless errands!" she
        vociferated. "Those are dames, who manage important matters and look after
        the house, and do you send them to ask for eatables and inquire about tea?
        You haven't even the least notion about gradation. P'ing Erh is standing
        here, so tell her to go and give the message." P'ing Erh immediately
        assented, and issued from the room, bent upon going on the errand. But the
        married women stealthily pulled her back.}
}
{
}

\Verse{68}
{
    \zA{}
}
{
    \eA{"How could you, miss, be made to go and tell them?" they smiled. "We've got
        some one here, who can do so!" So saying, they dusted one of the stone steps
        with their handkerchiefs. "You've been standing so long," they observed,
        "that you must feel quite tired. Do sit in this sunny place and have a
        little rest." P'ing Erh took a seat on the step. Two matrons attached to the
        tea-room then fetched a rug and spread it out for her. "It's cold on those
        stones," they ventured; "this is, as clean as it can be. So, miss, do make
        the best of it, and use it!" P'ing Erh hastily forced a smile. "Many
        thanks," she replied. Another matron next brought her a cup of fine new tea.
        "This isn't the tea we ordinarily drink," she quietly smiled.}
}
{
}

\Verse{69}
{
    \zA{}
}
{
    \eA{"This is really for entertaining the young ladies with. Miss, pray moisten
        your mouth with some." P'ing Erh lost no time in bending her body forward
        and taking the cup. Then pointing at the company of married women, she
        observed in a low voice: "You're all too fond of trouble! The way you're
        going on won't do at all! She (T'an Ch'un) is only a young girl, so she is
        loth to show any severity, or display any temper. This is because she's full
        of respect. Yet you people look down on her and insult her. Should she,
        however, be actually provoked into any violent fit of anger, people will
        simply say that her behaviour was rather rough, and all will be over. But as
        for you, you'll get at once into endless trouble.}
}
{
}

\Verse{70}
{
    \zA{}
}
{
    \eA{Even though she might show herself somewhat wilful, Madame Wang treats her
        with considerable forbearance, and lady Secunda too hasn't the courage to
        meddle with her; and do you people have such arrogance as to look down on
        her? This is certainly just as if an egg were to go and bang itself against
        a stone!" "When were we ever so audacious?" the servants exclaimed with one
        voice. "This fuss is all the work of Mrs. Chao!" "Never mind about that!"
        P'ing Erh urged again in an undertone. "My dear ladies, 'when a wall falls,
        every one gives it a shove.' That Mrs. Chao has always been rather
        topsy-turvey in her ways, and done things by halves; so whenever there has
        been any rumpus, you've invariably shoved the blame on to her shoulders.
        Never have you had any regard for any single person.}
}
{
}

\Verse{71}
{
    \zA{}
}
{
    \eA{Your designs are simply awful! Is it likely that all these years that I've
        been here, I haven't come to know of them? Had our lady Secunda mismanaged
        things just a little bit, she would have long ago been run down by every one
        of you, ladies! Even such as she is, you would, could you only get the least
        opportunity, be ready to place her in a fix! And how many, many times hasn't
        she been abused by you?" "She's dreadful," one and all of them rejoined.
        "You all live in fear and trembling of her. But we know well enough that no
        one could say that she too does not in the depths of her heart entertain
        some little dread for the lot of you. The other day, we said, in talking
        matters over, that things could not go on smoothly from beginning to end,
        and that some unpleasantness was bound to happen.}
}
{
}

\Verse{72}
{
    \zA{}
}
{
    \eA{Miss Tertia is, it's true, a mere girl, and you've always treated her with
        little consideration, but out of that company of senior and junior young
        ladies, she is the only soul whom our lady Secunda funks to some certain
        extent. And yet you people now won't look up to her." So speaking Ch'iu Wen
        appeared to view. The married women ran up to her and inquired after her
        health. "Miss," they said, "do rest a little. They've had their meal served
        in there, so wait until things have been cleared away, before you go and
        deliver your message." "I'm not like you people," Ch'iu Wen smiled. "How can
        I afford to wait?" With these words on her lips, she was about to go into
        the hall, when P'ing Erh quickly called her back. Ch'iu Wen, upon turning
        her head round, caught sight of P'ing Erh.}
}
{
}

\Verse{73}
{
    \zA{}
}
{
    \eA{"Have you too," she remarked with a smile, "come here to become something
        like those guardians posted outside the enclosing walls?" Retracing, at the
        same time, her footsteps, she took a seat on the rug, occupied by P'ing Erh.
        "What message have you got to deliver?" P'ing Erh gently asked. "I've got to
        ask when we can get Pao-yü's monthly allowance and our own too," she
        responded. "Is this any such pressing matter?" P'ing Erh answered. "Go back
        quick, and tell Hsi Jen that my advice is that no concern whatever should be
        brought to their notice to-day. That every single matter reported is bound
        to be objected to; and that even a hundred will just as surely be vetoed."
        "Why is it?" vehemently inquired Ch'iu Wen, upon hearing this explanation.
        P'ing Erh and the other servants then promptly told her the various reasons.}
}
{
}

\Verse{74}
{
    \zA{}
}
{
    \eA{"She's just bent," they proceeded, "upon finding a few weighty concerns in
        order to establish, at the expense of any decent person who might chance to
        present herself, a precedent of some kind or other so as to fix upon a mode
        of action, which might help to put down expenses to their proper level, and
        afford a lesson to the whole household; and why are you people the first to
        come and bump your heads against the nails? If you went now and told them
        your errand, it would also reflect discredit upon our venerable old mistress
        and Madame Wang, were they to pounce upon one or two matters to make an
        example of you.}
}
{
}

\Verse{75}
{
    \zA{}
}
{
    \eA{But if they complied with one or two of your applications, others will again
        maintain 'that they are inclined to favour this one and show partiality to
        that one; that as you had your old mistress' and Madame Wang's authority to
        fall back upon, they were afraid and did not presume to provoke their
        displeasure; that they only avail themselves of soft-natured persons to make
        scapegoats of.' Just mark my words! She even means to raise objections in
        one or two matters connected with our lady Secunda, in order to be the
        better able to shut up people's mouths." Ch'iu Wen listened to her with
        patient ear; and then stretching out her tongue, "It's lucky enough you were
        here, sister P'ing," she smiled; "otherwise, I would have had my nose well
        rubbed on the ground.}
}
{
}

\Verse{76}
{
    \zA{}
}
{
    \eA{I shall seize the earliest opportunity and give the lot of them a hint."
        While replying, she immediately rose to her feet and took leave of them.
        Soon after her departure, Pao-ch'ai's eatables arrived, and P'ing Erh
        hastened to enter and wait on her. By that time Mrs. Chao had left, so the
        three girls seated themselves on the wooden bed, and went through their
        repast. Pao-ch'ai faced the south. T'an Ch'un the west. Li Wan the east. The
        company of married women stood quietly under the verandah ready to answer
        any calls. Within the precincts of the chamber, only such maids remained in
        waiting as had ever been their closest attendants. None of the other
        servants ventured, of their own accord, to put their foot anywhere inside.
        The married women (meanwhile) discussed matters in a confidential whisper.}
}
{
}

\Verse{77}
{
    \zA{}
}
{
    \eA{"Let's do our downright best to save trouble," they argued. "Don't let us
        therefore harbour any evil design, for even dame Wu will, in that case, be
        placed in an awkward fix. And can we boast of any grand honours to expect to
        fare any better?" While they stood on one side, and held counsel together,
        waiting for the meal to be over to make their several reports, they could
        not catch so much as the caw of a crow inside the rooms. Neither did the
        clatter of bowls and chopsticks reach their ears. But presently, they
        discerned a maid raise the frame of the portiere as high as she could, and
        two other girls bring the table out.}
}
{
}

\Verse{78}
{
    \zA{}
}
{
    \eA{In the tea-room, three maids waited with three basins in hand. The moment
        they saw the dining-table brought out, all three walked in. But after a
        brief interval, they egressed with the basins and rinsing cups. Shih Shu, Su
        Yün and Ying Erh thereupon entered with three covered cups of tea, placed in
        trays. Shortly however these three girls also made their exit. Shih Shu then
        recommended a young maid to be careful and attend to the wants (of their
        mistresses). "When we've had our rice," she added, "we'll come and relieve
        you. But don't go stealthily again and sit down!" The married women at
        length delivered their reports in a quiet and orderly manner; and as they
        did not presume to be as contemptuous and offhandish as they had been
        before, T'an Ch'un eventually cooled down.}
}
{
}

\Verse{79}
{
    \zA{}
}
{
    \eA{"I've got something of moment," she then observed to P'ing Erh, "about which
        I would like to consult your mistress. Happily, I remembered it just now, so
        come back as soon as you've had your meal. Miss Pao-ch'ai is also here at
        present, so, after we four have deliberated together, you can carefully ask
        your lady whether action is to be taken accordingly or not." P'ing Erh
        acquiesced and returned to her quarters. "How is it," inquired lady Feng,
        "that you've been away such an age?" P'ing Erh smiled and gave her a full
        account of what had recently transpired. "What a fine, splendid girl Miss
        Tertia is!" she laughingly ejaculated. "What I said was quite right! The
        only pity is that she should have had such a miserable lot as not to have
        been born of a primary wife." "My lady, you're also talking a lot of trash!"}
}
{
}

\Verse{80}
{
    \zA{}
}
{
    \eA{P'ing Erh smiled. "She, mayn't be Madame Wang's child, but is it likely that
        any one would be so bold as to point the finger of scorn at her, and not
        treat her like the others?" Lady Feng sighed. "How could you know
        everything?" she remarked. "She is, of course, the offspring of a concubine,
        but as a mere girl, she can't be placed on the same footing as a man! By and
        bye, when any one aspires to her hand, the sort of supercilious parties, who
        now tread the world, will, as a first step, ask whether this young lady is
        the child of a No. 1 or No. 2 wife. And many of these won't have anything to
        say to her, as she is the child, of a No. 2.}
}
{
}

\Verse{81}
{
    \zA{}
}
{
    \eA{But really people haven't any idea that, not to speak of her as the
        offspring of a secondary wife, she would be, even as a mere servant-girl of
        ours, far superior than the very legitimate daughter of any family. Who, I
        wonder, will in the future be so devoid of good fortune as to break off the
        match; just because he may be inclined to pick and choose between a wife's
        child and a concubine's child? And who, I would like to know, will be that
        lucky fellow, who'll snatch her off without any regard to No. 1 and No. 2?"
        Continuing, she resumed, turning smilingly towards P'ing Erh, "You know well
        enough how many ways and means I've had all these years to devise in order
        to effect retrenchment, and how there isn't, I may safely aver, a single
        soul in the whole household, who doesn't detest me behind my back.}
}
{
}

\Verse{82}
{
    \zA{}
}
{
    \eA{But now that I'm astride on the tiger's back, (I must go on; for if I put my
        foot on the ground, I shall be devoured). It's true, my tactics have been
        more or less seen through, but there's no help for it; I can't very well
        become more open-handed in a moment! In the second place, much goes out at
        home, and little comes in; and the hundred and one, large and small, things,
        which turn up, are still managed with that munificence so characteristic of
        our old ancestors. But the funds, that come in throughout the year, fall
        short of the immense sums of past days. And if I try again to effect any
        savings people will laugh at me, our venerable senior and Madame Wang suffer
        wrongs, and the servants abhor me for my stinginess.}
}
{
}

\Verse{83}
{
    \zA{}
}
{
    \eA{Yet, if we don't seize the first opportunity to think of some plan for
        enforcing retrenchment, our means will, in the course of a few more years,
        be completely exhausted." "Quite so!" assented P'ing Erh. "By and bye, there
        will be three or four daughters and two or three more sons added; and our
        old mistress won't be able, singlehanded, to meet all this heavy outlay." "I
        myself entertain fears on the same score," lady Feng smiled. "But, after
        all, there will be ample. For when Pao-yü and cousin Lin get married, there
        won't be any need to touch a cent of public money, as our old lady has her
        own private means, and she can well fork out some. Miss Secunda is the child
        of your senior master yonder, and she too needn't be taken into account.}
}
{
}

\Verse{84}
{
    \zA{}
}
{
    \eA{So there only remain three or four, for each of whom one need only spend, at
        the utmost, ten thousand taels. Cousin Huan will marry in the near future;
        and if an outlay of three thousand taels prove insufficient, we will be
        able, by curtailing the bandoline, used in those rooms for smoothing the
        hair with, make both ends meet. And should our worthy senior's end come
        about, provision for everything is already made. All that we'll have to do
        will be to spend some small sum for a few miscellaneous trifles; and three
        to five thousand taels will more than suffice. So with further economies at
        present, there will be plenty for all our successive needs. The only fear is
        lest anything occur at an unforeseen juncture; for then it will be dreadful!
        But don't let us give way to apprehensions with regard to the future!}
}
{
}

\Verse{85}
{
    \zA{}
}
{
    \eA{You'd better have your rice; and when you've done, be quick and go and hear
        what they mean to treat about in their deliberations. I must now turn this
        opportunity to the best account. I was only this very minute lamenting that
        I had no help at my disposal. There's Pao-yü, it's true, but he too is made
        of the same stuff as the rest of them in here. Were I even to get him under
        my thumb, it would be of no earthly use whatever. Senior lady is as
        good-natured as a joss; and she likewise is no good. Miss Secunda is worse
        than useless. Besides, she doesn't belong to this place. Miss Quarta is only
        a child. That young fellow Lan and Huan-erh are, more than any of the
        others, like frozen kittens with frizzled coats.}
}
{
}

\Verse{86}
{
    \zA{}
}
{
    \eA{They only wait to find some warm hole in a stove into which they may poke
        themselves! Really from one and the same womb have been created two human
        beings (T'an Ch'un and Chia Huan) so totally unlike each other as the
        heavens are distant from the earth. But when I think of all this, I feel
        quite angry! Again, that girl Lin and Miss Pao are both deserving enough,
        but as they also happen to be our connexions, they couldn't very well be put
        in charge of our family affairs. What's more, the one resembles a lantern,
        decorated with nice girls, apt to spoil so soon as it is blown by a puff of
        wind. The other has made up her mind not to open her month in anything that
        doesn't concern her.}
}
{
}

\Verse{87}
{
    \zA{}
}
{
    \eA{When she's questioned about anything, she simply shakes her head, and
        repeats thrice: 'I don't know,' so that it would be an extremely difficult
        job to go and ask her to lend a helping hand. There's only therefore Miss
        Tertia, who is as sharp of mind as of tongue. She's besides a
        straightforward creature in this household of ours and Madame Wang is
        attached to her as well. It's true that she outwardly makes no display of
        her feelings for her, but it's all that old thing Mrs. Chao, who has done
        the mischief, for, in her heart, she actually holds her as dear as she does
        Pao-yü. She's such a contrast to Huan-erh! He truly makes it hard for any
        one to care a rap for him. Could I have had my own way, I would long ere
        this have packed him out of the place.}
}
{
}

\Verse{88}
{
    \zA{}
}
{
    \eA{But since she (T'au Ch'un) has now got this idea into her mind, we must
        cooperate with her. For if we can afford each other a helping hand, I too
        won't be single-handed and alone. And as far as every right principle,
        eternal principle, and honesty of purpose go, we shall with such a person as
        a helpmate, be able to save ourselves considerable anxiety, and Madame
        Wang's interests will, on the other hand, derive every advantage. But, as
        far as unfairness and bad faith go, I've run the show with too malicious a
        hand, and I must turn tail and draw back from my old ways.}
}
{
}

\Verse{89}
{
    \zA{}
}
{
    \eA{When I review what I've done, I find that if I still push my tyrannical rule
        to the bitter end, people will hate me most relentlessly; so much so, that
        under their smiles they'll harbour daggers, and much though we two may then
        be able to boast of having four eyes and two heads between us, they'll
        compass our ruin, when they can at any moment find us off our guard. We
        should therefore make the best of this crisis, so that as soon as she takes
        the initiative and sets things in order, all that tribe of people may for a
        time lose sight of the bitter feelings they cherish against us, for the way
        we've dealt with them in the past. But there's another thing besides.}
}
{
}

\Verse{90}
{
    \zA{}
}
{
    \eA{I naturally know the great talents you possess, but I feel mistrust lest you
        should, by your own wits, not be able to bring things round. I enjoin these
        things then on you, now, for although a mere girl she has everything at her
        fingers' ends. The only thing is that she must try and be wary in speech.
        She's besides so much better read than I am that she's a harder nut to
        crack. Now the proverb says: 'in order to be able to catch the rebels, you
        must first catch their chief.' So if she's at present disposed to mature
        some plan and set to work to put it into practice, she'll certainly have to
        first and foremost make a start with me. In the event consequently of her
        raising objections to anything I've done, mind you don't begin any dispute
        with her.}
}
{
}

\Verse{91}
{
    \zA{}
}
{
    \eA{The more virulent she is in her censure of me, the more deferential you
        should be towards her. That's your best plan. And whatever you do, don't
        imagine that I'm afraid of any loss of face. But the moment you flare up
        with her, things won' go well……" P'ing Erh did not allow her time to
        conclude her argument. "You're too much disposed to treat us as simpletons!"
        she smiled. "I've already carried out your wishes, and do you now enjoin all
        these things on me?" Lady Feng smiled. "It's because," she resumed, "I
        feared lest you, who have your eyes and mouth so full of me, and only me,
        might be inclined to show no regard whatever for her, that's why. I
        couldn't, therefore, but tender you the advice I did.}
}
{
}

\Verse{92}
{
    \zA{}
}
{
    \eA{But since you've already done what I wanted you to do, you've shown yourself
        far sharper than I am. There's nothing in this to drive you into another
        tantrum, and to make that mouth of yours begin to chatter away so much about
        'you and I,' 'you and I' !" "I've actually addressed you as 'you' ;" P'ing
        Erh rejoined; "but if you be displeased at it, isn't this a case of a slap
        on the mouth? You can very well give me another one, for is it likely that
        this phiz of mine hasn't as yet tasted any, pray?" "What a vixen you are!"
        lady Feng said smilingly. "How many faults will you go on picking out,
        before you shut up? You see how ill I am, and yet you come to rub me the
        wrong way. Come and sit down; for you and I can at all events have our meal
        together when there is no one to break in upon us.}
}
{
}

\Verse{93}
{
    \zA{}
}
{
    \eA{It's only right that we should." While these remarks dropped from her lips,
        Feng Erh and some three or four other maids entered the room and laid the
        small stove-couch table. Lady Feng only ate some birds' nests' soup and
        emptied two small plates of some recherché light viands; for she had long
        ago temporarily reduced her customary diet. Feng Erh placed the four kinds
        of eatables allotted to P'ing Erh on the table. After which, she filled a
        bowl of rice for her. Then with one leg bent on the edge of the stove-couch,
        while the other rested on the ground, P'ing Erh kept lady Feng company
        during her repast; and waiting on her, afterwards, until she finished
        rinsing her mouth, she issued certain directions to Feng Erh, and crossed
        over at length to T'an Ch'un's quarters.}
}
{
}

\Verse{94}
{
    \zA{}
}
{
    \eA{Here she found the courtyard plunged in perfect stillness, for the various
        inmates, who had been assembled there, had already taken their leave. But,
        reader, do you wish to follow up the story? If so, listen to the
        circumstances detailed in the next chapter.}
}
{
}
